%======================================================================
\sloppy
\section{Experimental setup}\label{sec:setup}
%======================================================================

\subsection{Datasets}\label{sec:datasets}

We use four public benchmarks. Each pairs a biased log with a test set collected without the
selection of the log, so that a debiasing method can be measured against unbiased outcomes.
Table~\ref{tab:data} gives the sizes, which we checked against the data loaders in the repository.

\begin{table}[t]\centering\scriptsize
\caption{Datasets. ``Log'': the biased training data, with its density. ``Test'': the unbiased
data. ``Candidates'': unbiased test items per user that are ranked (items never logged for that
user). ``Evaluated users'': users with at least one positive candidate in the full test set.}\label{tab:data}
\begin{adjustbox}{max width=\textwidth}
\begin{tabular}{lcccccc}
\toprule
Dataset & Users / items & Log (density) & $O$ means & Test & Candidates & Evaluated users\\
\midrule
Coat & 290 / 300 & 6{,}960 (8.0\%) & Self-rated & 16 random ratings per user & 14.7 & 225\\
Yahoo!\,R3 & 15{,}400 / 1{,}000 & 311{,}704 (2.0\%) & Self-rated & 10 random songs, 5{,}400 users & 10 & 2{,}315\\
KuaiRand & 10{,}000 / 7{,}583 & 511{,}863 (0.68\%) & Feed impression & random impressions & 43.6 & 8{,}347\\
KuaiRec & 7{,}176 / 10{,}728 & 10.3M (13.4\%) & Impression & fully observed 1{,}411 $\times$ 3{,}327 & 3{,}314 & 1{,}411\\
\bottomrule
\end{tabular}

\end{adjustbox}
\end{table}

\emph{Coat}~\citep{schnabel2016} has 290 users and 300 items. The log holds 6{,}960
self-selected ratings (8.0\% of the matrix), of which 1{,}905 are ratings of 4 or more. The test
set holds 16 ratings per user on uniformly random items. Candidates are the items a user did not
rate in the log, 14.7 per user on average. The dataset ships propensities, which its authors estimated from a uniform sample; the
propensity-based methods use them on Coat (Section~\ref{sec:nuisance}).
\emph{Yahoo!\,R3}~\citep{marlin2009collaborative} has 15{,}400 users and 1{,}000 songs. The log holds
311{,}704 self-selected ratings (2.0\% dense). The test set holds 54{,}000 ratings of ten random
songs for 5{,}400 users, in the version distributed with~\citep{chen2021autodebias}. A random 5\%
of these users (270, fixed seed) is excluded from validation and testing. The remaining 5{,}130
test users have ten candidates each, and 2{,}315 of them have a positive candidate.
\emph{KuaiRand-Pure}~\citep{gao2022kuairand} is the only benchmark whose log records platform
exposure: an impression is an item shown in the standard feed, and the outcome is a click. We keep
10{,}000 users who also received randomly exposed videos, and all 7{,}583 videos. The log holds
511{,}863 impressions (0.68\% dense), 47\% of them clicks. The test set consists of the random
impressions of the same period, restricted to videos the user was not shown in the standard log.
Users have 43.6 candidates on average, and 8{,}347 of them have a positive candidate.
\emph{KuaiRec}~\citep{gao2022kuairec} provides a dense log of 10.3 million views by 7{,}176 users of
10{,}728 items (13.4\% dense). Its test set is a fully observed matrix for 1{,}411 users and 3{,}327
items. Each user is ranked over the test items that the user never viewed in the log, 3{,}314 on
average.

Positive outcomes are ratings of 4 or more (Coat, Yahoo!\,R3), a click (KuaiRand) and a mean watch
ratio of at least 2 (KuaiRec). On Coat and Yahoo!\,R3, ``exposure'' means that the user chose to
rate the item, so the propensity models self-selection and not a platform's allocation. The
candidate sets of these two datasets hold about ten to fifteen items, which limits the range that
nDCG can span on them. The ``evaluated users'' column of Table~\ref{tab:data} counts every user with
a positive candidate. The test-user counts in Section~\ref{sec:protocol} are slightly lower for
Yahoo!\,R3, KuaiRand and KuaiRec, because a user can fall in the validation part of every split and
then has no test entry.

\subsection{Baselines}\label{sec:baselines}

Every method is run by us on the same splits and data preparation. No number is taken from the
literature. Graph models are LightGCN~\citep{he2020lightgcn}, r-AdjNorm~\citep{zhao2022radjnorm},
NAVIP~\citep{kim2022navip}, SimGCL~\citep{yu2022simgcl}, DR-LightGCN and LightGCN-pt. Matrix models
are iALS~\citep{hu2008ials,rendle2022ials}, BPR-MF, MACR~\citep{wei2021macr}, DR-JL~\citep{wang2019drjl} and
MRDR~\citep{guo2021mrdr}. AdvInfoNCE is the contrastive baseline (Section~\ref{sec:advsetup}).
The training-free operators are DRUP and DR adjacency, which is DRUP without the walk correction.
On Coat, the propensity-based baselines (DR-LightGCN, NAVIP, DR-JL, MRDR, DR adjacency and DRUP) use the
same shipped propensities as the method.

LightGCN is trained with BPR on the logged positives and uniform negatives. It is the loss of
Section~\ref{sec:weighted} with unit weights and no offset. LightGCN-pt is LightGCN trained with a
pointwise cross-entropy on all logged pairs (KuaiRec only). r-AdjNorm searches the normalisation
exponent $r\in\{0.3,0.7\}$ on top of the LightGCN grid. NAVIP trains LightGCN with BPR on a
neighbourhood aggregation in which each logged edge is weighted by the inverse of its estimated
propensity. DR-LightGCN uses the pointwise doubly robust loss with the propensity and imputation of
\qp, and pairs drawn uniformly from all pairs. SimGCL adds uniform noise to the propagated
embeddings and an InfoNCE term between two noisy views, with weight in $\{0.1,0.5\}$. iALS puts
confidence $1+\alpha$ on logged positives and 1 on the other pairs, with $\alpha\in\{3,10,30\}$ and
regularisation in $\{1,10,100,1000\}$. BPR-MF is matrix factorisation trained with BPR. Our MACR
baseline is a MACR-style model on a matrix factorisation backbone, with user and item branches and
popularity removed at inference ($c\in\{0.1,0.3\}$); it is not the original code. DR-JL and MRDR fit
matrix factorisation with the joint imputation of their losses.

The final comparison contains the following methods. Coat: LightGCN, r-AdjNorm, NAVIP, DR-LightGCN,
iALS, SimGCL, MACR and AdvInfoNCE. Yahoo!\,R3: LightGCN, r-AdjNorm, NAVIP, iALS, SimGCL and
AdvInfoNCE. KuaiRand: LightGCN, r-AdjNorm, NAVIP, iALS, BPR-MF, SimGCL and AdvInfoNCE. KuaiRec:
DR-JL, MRDR, LightGCN-pt, DR-LightGCN, NAVIP, iALS, LightGCN, r-AdjNorm and AdvInfoNCE. These are
the Holm families of Section~\ref{sec:protocol}. DRUP, DR adjacency, and DR-JL and MRDR on Coat and
Yahoo!\,R3 are reported without a test, because they are outside the families.

We did not run StableDR~\citep{li2023stabledr}, DICE~\citep{zheng2021dice}, CausE~\citep{bonner2018cause}
or AutoDebias~\citep{chen2021autodebias}. CausE and AutoDebias need uniformly collected data in
training, which our protocol reserves for validation and testing. For StableDR and DICE, the DR and
popularity-aware baselines we ran are the closest counterparts. We did not compare with LLM-based
recommenders. None of the four benchmarks provides item text together with a random-exposure test
set, so such a comparison would test a different hypothesis.

\subsection{AdvInfoNCE re-implementation}\label{sec:advsetup}

AdvInfoNCE~\citep{zhang2023advinfonce} is an adversarial contrastive loss. We re-implemented it on
the same LightGCN encoder, and our numbers are for this re-implementation, not for the original
paper. The main simplifications are these.

\begin{itemize}
\item The score is the cosine similarity of the propagated embeddings divided by a temperature
$\tau$. Each positive is contrasted with 64 sampled negatives.
\item The adversary is a pair of embedding tables (users and items, dimension 64) trained with Adam
to maximise the contrastive loss, with learning rate $10^{-2}$ on Coat, $10^{-3}$ on Yahoo!\,R3 and
KuaiRand, and $5\cdot10^{-4}$ on KuaiRec. Only this embedding-based adversary is implemented.
\item The adversary budget of the original formulation is replaced by a number of adversarial rounds
$\eta\in\{4,12\}$. A round runs after each epoch from the second epoch on and lasts at least one
pass over the positives and 100 Adam steps. The adversary's KL divergence is logged but not
projected onto a ball.
\end{itemize}

The grids differ in size. Coat has 216 configurations: three learning rates, three weight decays,
widths $\{32,64\}$, depths $\{1,2,3\}$, $\tau\in\{0.2,0.75\}$ and $\eta\in\{4,12\}$. Yahoo!\,R3 has 16
configurations with $\tau\in\{0.1,0.28\}$, $\eta\in\{4,12\}$, learning rates $\{3\cdot10^{-3},10^{-3}\}$,
weight decays $\{10^{-5},10^{-4}\}$, width 128 and depth 2. KuaiRand and KuaiRec have 8 each: KuaiRand
uses $\tau\in\{0.1,0.28\}$, learning rates $\{10^{-2},3\cdot10^{-3}\}$, weight decay $10^{-4}$, width 64
and depth 2; KuaiRec uses $\tau\in\{0.2,2.0\}$, learning rates $\{10^{-2},3\cdot10^{-3}\}$, weight decay
$10^{-5}$, width 64 and depth 2; both use $\eta\in\{4,12\}$. A larger grid gives a method more chances to
fit noise in the validation data, so the Coat grid favours AdvInfoNCE and the smaller grids work
against it.

\subsection{Protocol}\label{sec:protocol}

\paragraph{Splits.}
The uniformly collected data are split into a validation part and a test part. On Coat the split is
made within each user: 30\% of each user's candidates go to validation and the rest to test. On
Yahoo!\,R3, KuaiRand and KuaiRec the split is made over users: 30\% of the test users go to
validation, and no user is in both parts of a split. Coat uses ten random splits (seeds 0 to 9) and
the other datasets use five (seeds 0 to 4).

\paragraph{Training and selection.}
Each configuration is trained once on the log, and the same run serves every split. After each epoch
the model is scored on the validation and test parts of every split. Each split keeps its own best
(configuration, epoch) by validation nDCG and stops when its validation score has not improved for
the patience period: five epochs on Coat and Yahoo!\,R3, three on KuaiRand and KuaiRec. A run ends
when every split has stopped or the epoch limit is reached. Test scores are stored for every epoch,
but they never influence the choice. All methods on a dataset share the epoch limit, batch size and
patience, except AdvInfoNCE on KuaiRec, which uses batch size 4{,}096 instead of 32{,}768.

\paragraph{Metrics.}
We report nDCG@5 on Coat and Yahoo!\,R3, nDCG@10 on KuaiRand and nDCG@20 on KuaiRec. Each user's
candidates are ranked within the part being scored, and users with no positive candidate in that part
are skipped. We rank all candidates instead of sampling negatives, because sampled metrics are not
consistent with the full ranking metric~\citep{krichene2020sampled}. We report the mean and the
standard deviation over splits (with $\mathrm{ddof}=0$). Evaluation uses only the uniformly collected
test data, following the recommendation of~\citet{steck2010mnar} to evaluate on data that are not
missing not at random.

\paragraph{Tests.}
Differences are tested at the level of users. For each user we average the user's test nDCG over the
splits in which the user is a test user. For a method and a baseline we take the per-user differences,
apply a one-sample $t$-test, and report the 95\% confidence interval as a half-width. Within each
dataset the $p$-values of the family are adjusted by the Holm procedure~\citep{holm1979}. A difference
is called significant only if the adjusted $p$ is below 0.05. Differences are reported as the mean
per-user difference times $10^{3}$, with the half-width in brackets. This is not the difference of the
two table means. The Holm families have 8 tests on Coat, 6 on Yahoo!\,R3, 7 on KuaiRand and 9 on
KuaiRec. The number of users with a test entry in the final method's files is 225, 2{,}310, 8{,}321 and
1{,}410, respectively.

Splits overlap, so per-user values are not independent across splits, and paired tests over
overlapping splits can be anti-conservative~\citep{nadeau2003}. The tests treat the training log, the
fitted nuisances and the trained models as fixed. We do not report a corrected resampled test for the
final method.

\subsection{Seed robustness}\label{sec:seedsetup}

Each split's selected configuration is retrained with seeds 1, 2 and 3. Each split re-selects its
epoch on its own validation part, and the configuration is not changed. The retrained final method is
compared with the two highest-scoring baselines of the main run for each dataset, where ``highest-scoring'' means
highest main-run test mean. The family of two is therefore chosen on the test data, and its tests are descriptive: iALS and DR-LightGCN
on Coat, r-AdjNorm and LightGCN on Yahoo!\,R3 and KuaiRand, and DR-JL and MRDR on KuaiRec. The Holm
correction runs over those two tests. The two baselines are retrained in the same way: for each split, the
configuration selected in their main-run files is retrained with seeds 1 to 3, and the epoch is re-selected on
validation data, as for the final method. Their seed-0 values are the main-run values. Some main-run files also contain a re-run of split 0 with seeds
100 to 102, stored in a separate field. These runs are not used in any table.

\subsection{Tuning grids and budgets}\label{sec:grids}

The final method is the union of several jobs, one per group of $\beta$ values and learning rate (and,
on Coat, one per value of $f_n$ in the offset-free job). Each split keeps the job with the best
validation score (Section~\ref{sec:offset}).

\begin{itemize}
\item \textbf{Coat.} Learning rate $\{10^{-2},3\cdot10^{-3},10^{-3}\}$, weight decay
$\{10^{-5},10^{-4},10^{-3}\}$, width $\{32,64\}$, depth $\{1,2,3\}$, at most 100 epochs and batch size
256. The three jobs (offset $\beta=0$ with $f_n\in\{0.5,1\}$; $\beta\in\{0.5,1\}$; $\beta\in\{1,2\}$)
have 108 configurations each, 324 in total.
\item \textbf{Yahoo!\,R3.} Learning rate $\{10^{-2},3\cdot10^{-3},10^{-3}\}$ for the $\beta=0$ jobs and
$3\cdot10^{-3}$ for the offset jobs, weight decay $\{10^{-5},10^{-4}\}$, width $\{64,128\}$, depth
$\{1,2,3\}$, at most 60 epochs and batch size 1{,}024. The five jobs have 84 configurations in total.
\item \textbf{KuaiRand.} Learning rate $\{10^{-2},3\cdot10^{-3},10^{-3}\}$, weight decay
$\{10^{-5},10^{-4}\}$, width 64, depth 2, at most 40 epochs and batch size 8{,}192. The three jobs have
30 configurations in total.
\item \textbf{KuaiRec.} Weight decay $\{10^{-5},10^{-4}\}$, width $\{64,128\}$, depth $\{1,2,3\}$, at most
30 epochs and batch size 32{,}768. Learning rates are $\{10^{-2},3\cdot10^{-3}\}$ for the $\beta=0$ jobs,
$\{10^{-2},3\cdot10^{-3},10^{-3}\}$ for $\beta\in\{1,2\}$, and $3\cdot10^{-3}$ for $\beta\in\{0.5,1\}$. The
six jobs have 120 configurations in total. The cap of 2{,}000{,}000 pairs per epoch never binds, because
the log holds 779{,}468 positives.
\end{itemize}

The baselines use the same learning-rate and weight-decay values and their own parameters from
Section~\ref{sec:baselines}. The number of configurations searched for each baseline is in the ``cfgs''
column of the main results table. The baseline counts range from 6 to 216 (AdvInfoNCE on Coat), and
the full grid of each baseline will be listed in the supplement. The unequal budgets are a limitation
of the comparison, since selection over more configurations can favour the method.
